\begin{helper}[A PE support appears by its supporting slot]
Let \(\mathcal T\) be a certified Tashkinov series and let \(\mathcal E\)
be an ETT carrying that series, with displayed tree \(T\).  If the positive
data slot \(q^+\) is PE with chosen supporting vertex \(v\), and
\(h:v\in V(T)\), then
\[
 \noderef{ETTFirstAppearanceIndex}(\mathcal E,v,h)\le q+1.
\]
Thus a supporting vertex cannot first occur only after the slot at which it
is used.
\end{helper}
\begin{proof}
The certificate at \(q^+\), supplied by \noderef{TashkinovSeries}, cannot
be initial because the index is a successor.  Its displayed PE-kind
equation rules out the SE and RE constructors of
\noderef{TashkinovSeriesStep}.  In the remaining PE constructor, equality
of the chosen-vertex fields identifies its
\noderef{MaximumDefectiveChoice} vertex with \(v\).  That vertex is
defective on the ladder tree at \(q^+\), so in particular belongs to that
tree's vertex set.

By \noderef{ExtendedTashkinovLadderStage}, the ladder stage with zero-based
index \(q\) is exactly this tree.  It is therefore a member of the finite
candidate set minimized in the definition
\noderef{ETTFirstAppearanceIndex}.  The minimum candidate is at most \(q\);
adding one gives the asserted bound.
\end{proof}
